problem:  
Let \(G\) be a connected, simply connected, **equiregular Carnot group** with graded Lie algebra  
\(\mathfrak g = V_1 \oplus \cdots \oplus V_s\) and homogeneous dimension \(Q\).  
Equip \(V_1\) with a left‑invariant sub‑Riemannian metric \(g_\Delta\) and let \(d_\Delta\) denote the corresponding Carnot–Carathéodory distance.  
For each integer \(1 \le k < \dim(G)\), define the \(k\)-dimensional **isoperimetric profile**  
\[
\mathrm{Iso}_{G,g_\Delta}^{(k)}(r)=
\inf\Bigl\{
\mu_{k+1}(T)\;:\;
\partial T=S,\;
\mu_k(S)\le r
\Bigr\},
\]
where \(S,T\) are \(k\)- and \((k+1)\)-dimensional metric integral currents in \((G,d_\Delta)\) (Ambrosio–Kirchheim),  
and \(\mu_j\) denotes the \(j\)-dimensional Hausdorff measure for \(d_\Delta\).

It is known in special Carnot groups (e.g. Heisenberg) that  
\(\mathrm{Iso}^{(k)}(r)\) admits polynomial asymptotic behavior  
\(\mathrm{Iso}^{(k)}(r)\simeq C_{k}(G,g_\Delta)\, r^{\alpha(k,G)}\)  
for some exponent \(\alpha(k,G)\).  
For general Carnot groups only inequalities of this type are established.

---

**Problem (Algebraic determination of asymptotic isoperimetric profiles):**  
For each \(k\), let  
\[
\mathcal{C}^{*}_k(G,g_\Delta)
:= \limsup_{r\to \infty} \frac{\mathrm{Iso}^{(k)}_{G,g_\Delta}(r)}{r^{\alpha(k,G)}}
\in [0,\infty].
\]

1. Is \(\mathcal{C}^{*}_k(G,g_\Delta)\) determined **up to a universal multiplicative constant** solely by the graded Lie algebra \(\mathfrak g\)?  
   Equivalently, do two Carnot groups \(G,H\) with isomorphic graded Lie algebras satisfy  
   \[
   \mathcal{C}^{*}_k(G,g_\Delta) \asymp \mathcal{C}^{*}_k(H,h_\Delta)
   \]
   for all admissible horizontally inner products \(g_\Delta,h_\Delta\) (with constants independent of metric choice)?  

2. If not, can one build a quantitative invariant \(J_k(\mathfrak g,g_\Delta)\) describing how metric deformations of \(V_1\) alter \(\mathcal{C}^{*}_k\)?  
   In particular, does \(\mathcal{C}^{*}_k\) vary smoothly or discontinuously under changes of the metric shape operator on \(V_1\)?  

---

Why is it a "good" problem:  

- **Profound insight:**  
  This problem probes whether **asymptotic filling behavior**—a global geometric‑measure invariant—is encoded purely by algebraic layerings of the Lie algebra or whether the analytic choice of horizontal metric genuinely enriches the geometry. A positive answer would establish a striking universality law for all Carnot geometries; a negative one would precisely quantify how analytic anisotropy escapes the algebraic framework.

- **Bridge between fields:**  
  It interlaces **sub‑Riemannian geometry** (Carnot–Carathéodory metrics and stratified Lie algebras) with **geometric measure theory** (isoperimetric profiles, integral currents) and **asymptotic geometry** (large‑scale invariants). Any advance would force development of analytic tools sensitive to both bracket algebra and large‑scale measure growth.

- **Extreme simplicity:**  
  The essence can be stated in one line:
  > “Is the asymptotic isoperimetric profile of a Carnot group determined solely by its Lie algebra?”  
  Despite its compactness, the question penetrates one of the most fundamental boundaries in modern geometry—the tension between algebraic and analytic determinants of quantitative shape.